\section{Reliability 与 Generalized AI Systems}

最后十分钟把乐观 demo 收束为系统问题：agents 会循环、难以测试、可能失去监督，也会在缺少反馈时产生 plan divergence。课堂随后借 LLM OS 和 generalized AI system 图，尝试把 chat、rules、task engine、router、apps、local index 与 reflection 放进统一架构。

\subsection{Key Issues：可靠性、循环、评测与停止能力}

Key Issues slide 列出 reliability、looping、testing/benchmarking，以及 reset/deployment/observability 一类尚未成熟的问题。最关键的工程态度是：agent failure 不是单一错误率，而是一个动作如何改变后续 observation 和 decision 的动态过程。

\lecturefigure{slide-60-key-agent-issues.jpg}{Autonomous Agents 的关键可靠性问题}{Stanford 课堂视频 00:50:30}

\begin{knowledgebox}{读图：四类问题需要不同证据}
Reliability 需要任务成功率与错误分类；looping 需要 step limit 与重复状态检测；testing 需要可重放环境和 hidden test；oversight 需要 trace、kill switch 与责任人。只用最终回答 benchmark 无法覆盖真实行动系统。
\end{knowledgebox}

\begin{importantbox}{Agent Evaluation 的分层指标}
\begin{tabularx}{\textwidth}{L{0.22\textwidth}X X}
\toprule
层级 & 示例指标 & 不能忽略的失败 \\
\midrule
Step & action validity、tool error & 点击错对象、参数越界 \\
Trajectory & progress、loop rate、budget & 反复执行、状态漂移 \\
Task & verified success、human intervention & 看似完成但 postcondition 错 \\
Safety & unauthorized action、recovery & 越权、泄露、不可逆损失 \\
\bottomrule
\end{tabularx}
\end{importantbox}

\teachervoice{讲者强调必须能 stop agent；如果系统持续循环或执行错误动作，人类需要 reset、oversight 或 kill mechanism。自治不是移除监督，而是把监督从逐步操作改成边界和异常控制。}

\subsection{Plan Divergence：没有反馈的计划会离开现实}

Plan Divergence slide 用黑色 ideal plan 与红色 agent trajectory 表示偏离。误差一开始可能很小，但若系统只根据自己生成的文本继续规划，不读取网页状态、compiler error 或用户修正，偏差会累积并形成新的错误前提。本节回到 feedback loop，说明怎样用外部可验证信号把 trajectory 拉回任务目标。

\lecturefigure{slide-61-plan-divergence.jpg}{缺乏环境反馈时的 Plan Divergence}{Stanford 课堂视频 00:53:09}

\begin{importantbox}{把 Agent 看成反馈控制系统}
令期望状态为 $x_t^*$，实际环境状态为 $x_t$，误差为 $e_t=x_t^*-x_t$。无反馈 open-loop 规划只执行预生成动作；closed-loop agent 在每步读取 observation，并用
\[
a_t=\pi(g,x_t,e_t),\qquad x_{t+1}=F(x_t,a_t)
\]
修正后续动作。关键不是公式复杂，而是 $x_t$ 必须来自环境证据，不能由模型想象。
\end{importantbox}

\begin{warningbox}{自我反思不是环境验证的替代品}
让同一个模型阅读自己的计划，可能改善明显矛盾，却无法知道网页是否提交、代码是否编译或库存是否变化。优先使用 compiler、tests、API response、DOM state、receipt 等 external verifier；语言 critique 只能作为补充。
\end{warningbox}

\teachervoice{课堂以 coding agent 为正例：compiler 或 IDE 能返回明确错误，使 agent 知道“你错了”。讲者把早期 AutoGPT-style 系统的 divergence 归因于缺少可靠 feedback；这比单纯增加思考轮次更接近根因。}

\subsection{LLM OS：围绕模型建立操作层}

Karpathy LLM OS slide 把 LLM 放在中央，周围连接 conversational interface、tools、browser、storage、vision 与其他 I/O。它不是一个已经标准化的操作系统，而是 architecture metaphor：模型像 compute kernel，系统层负责上下文、资源、路由和安全。

\lecturefigure{slide-62-llm-os.jpg}{Karpathy LLM OS 的课堂架构图}{Stanford 课堂视频 00:53:33}

\begin{knowledgebox}{读图：OS 类比对应哪些职责}
Chat interface 类似 user shell；context/memory 类似 working set 与 storage；router 类似 scheduler；tools/apps 类似 devices 与 services；permissions/sandbox 类似 protection boundary；reflection/verifier 类似 error handling。类比用于发现缺失组件，不代表 LLM 具备传统 OS 的确定性与隔离保证。
\end{knowledgebox}

\teachervoice{讲者说，人类没有一直把 CPU 当作“电脑本身”，同样也不应一直把 chatbot 当作 AI 的最终形态。系统需要把 model 包在规则、任务引擎、存储和工具之外，才能承担更复杂工作。}

\subsection{Generalized AI System：chat、rules、routing 与 reflection}

Generalized AI Systems slide 展示更完整的数据流：用户通过 chat interface 输入，系统结合 rules 创建 task，经 task engine 与 router 分发到 apps/tools 或 local index，再把结果送回 reflection 与用户。读图时应从控制顺序而不是组件名称出发。

\lecturefigure{slide-63-generalized-ai-systems.jpg}{Building Generalized AI Systems 的课堂草图}{Stanford 课堂视频 00:58:51}

\begin{knowledgebox}{读图：一条安全任务路径}
输入先经过规则检查和目标澄清，task engine 生成 bounded plan，router 选择 browser、app、tool 或 local index；各执行器返回 evidence，reflection 检查是否满足 success predicate。若不满足，系统在预算内修订；若触及权限或风险边界，则交还用户。
\end{knowledgebox}

\begin{warningbox}{Rules 必须在执行前约束计划}
若禁止事项只写在最终回复模板里，agent 仍可能已经完成越权动作。规则要参与 plan validation、tool authorization 和 argument checking；任何违反不可变原则的候选计划都应在执行前被拒绝。
\end{warningbox}

\teachervoice{课堂提出“inherent rules”：系统生成的计划若违反原则，应自动 invalidated。随后 manager/router 决定访问 browser、apps、tools 或 local files，reflection 再判断任务是否正确完成。}

\subsection{Future Needs：error correction、permissions 与 sandboxing}

最后一页列出 error correction、security/user permissions、sandboxing 与高风险部署。它把整堂 Agent 讨论落到最现实的验收条件：不仅要能完成顺利路径，还要能捕获错误、限制访问、隔离执行，并保护 business、financial 和其他不可逆操作。

\lecturefigure{slide-64-future-needs.jpg}{未来 Agent 需要的纠错、安全与隔离机制}{Stanford 课堂视频 01:00:06}

\begin{importantbox}{部署前的最小安全闭环}
\begin{enumerate}
  \item \textbf{Least privilege}：每个 tool 只获得完成当前任务所需权限；
  \item \textbf{Sandbox}：代码、文件和浏览操作在受限环境执行；
  \item \textbf{Confirmation gate}：付款、删除、发布、发送和账号变更前确认；
  \item \textbf{Audit trail}：保存目标、计划、动作、观察、授权与结果；
  \item \textbf{Recovery}：支持取消、补偿事务、版本回滚和人工接管；
  \item \textbf{Evaluation}：在 adversarial 与长时任务中测试越权、注入和 divergence。
\end{enumerate}
\end{importantbox}

\teachervoice{讲者最后用非常具体的例子解释 permissions：agent 可以访问 DoorDash，但未必能访问 bank account；还要防止删除电脑内容，并对商业、金融和不可逆风险设置保护。这些不是未来附加功能，而是可部署 agent 的组成部分。}

\subsection{本章小结}

Agent reliability 必须在 trajectory 与 environment loop 上评测。Plan divergence 说明 external feedback 的必要性；LLM OS 与 generalized-system 图说明 chat、rules、task engine、router、tools、memory、reflection 和 permissions 应被设计为明确组件；future-needs 页则把 error correction、sandboxing 和 irreversible-action safeguards 设为部署门槛。

